\documentclass[11pt,a4paper]{article}

\usepackage[T1]{fontenc}
\usepackage[utf8]{inputenc}
\usepackage{lmodern}
\usepackage[italian]{babel}
\usepackage[margin=2.4cm]{geometry}
\usepackage{amsmath,amssymb}
\usepackage{microtype}
\usepackage{booktabs}
\usepackage{array}
\usepackage{xcolor}
\usepackage{caption}
\usepackage{listings}
\usepackage{tikz}
\usetikzlibrary{arrows.meta,positioning,calc,decorations.pathreplacing,shapes.geometric}
\usepackage[hidelinks]{hyperref}

\definecolor{accent}{HTML}{1F9D55}
\definecolor{hand}{HTML}{E0262B}
\definecolor{dot}{HTML}{22C443}

\lstdefinelanguage{GLSL}{
  language=C,
  morekeywords={vec2,vec3,vec4,uniform,sampler2D,texture2D,mix,smoothstep,clamp,max,min,length},
}
\lstset{
  language=GLSL,
  basicstyle=\ttfamily\small,
  keywordstyle=\color{accent!70!black}\bfseries,
  commentstyle=\color{gray}\itshape,
  frame=single,
  rulecolor=\color{gray!40},
  xleftmargin=1em,
  framexleftmargin=0.6em,
  columns=fullflexible,
}

\newcommand{\Sm}{\operatorname{S}}
\newcommand{\clamp}{\operatorname{clamp}}
\newcommand{\vect}[1]{\mathbf{#1}}

\title{La computer vision dietro \emph{Rectangular Filter}\\[4pt]
\large Mani, rettangoli, filtri e come far sparire una persona}
\author{Documentazione tecnica del progetto}
\date{Novembre 2026}

\begin{document}
\maketitle

\begin{abstract}
\noindent \emph{Rectangular Filter} è una web app che, a partire dal flusso della fotocamera,
riconosce le mani, costruisce un quadrilatero con le punte di pollici e indici e applica un
filtro solo al suo interno. Un gesto (pugno seguito da una V) fa sparire la persona dalla scena,
lasciando visibile solo lo sfondo. Questo documento spiega le tecniche di computer vision
e di elaborazione delle immagini usate a ogni passo: le reti neurali di MediaPipe per mani e
sagoma, la geometria del quadrilatero, il filtro One Euro, le funzioni di distanza con segno, i
filtri come operatori sull'immagine, la classificazione delle pose della mano e il compositing
con sottrazione dello sfondo. Tutti i numeri citati sono quelli effettivamente usati nel codice.
\end{abstract}

\setcounter{tocdepth}{2}
{\small\tableofcontents}
\clearpage

%=====================================================================
\section{La pipeline in un colpo d'occhio}
%=====================================================================

Ogni volta che la fotocamera produce un nuovo fotogramma (tipicamente 30 al secondo), l'app
esegue la catena di figura~\ref{fig:pipeline}. Tutto avviene sul dispositivo: le reti neurali
girano nel browser (WebAssembly o WebGL) e il compositing finale è uno \emph{shader} eseguito
dalla GPU, un pixel alla volta e in parallelo.

\begin{figure}[htbp]
\centering
\begin{tikzpicture}[
  box/.style={draw, rounded corners=3pt, align=center, minimum height=11mm, text width=25mm, font=\small},
  arr/.style={-{Stealth[length=2.2mm]}, thick, gray!70!black},
  lab/.style={font=\scriptsize, inner sep=1.5pt}]
  \node[box, fill=gray!10]   (cam)  at (0,0)      {Fotocamera\\\footnotesize fotogramma $I_t$};
  \node[box, fill=blue!8]    (hand) at (3.2,1.5)  {Hand\\Landmarker\\\footnotesize 21 punti per mano};
  \node[box, fill=purple!8]  (seg)  at (3.2,-1.5) {Selfie\\Segmenter\\\footnotesize maschera $S_t$};
  \node[box, fill=green!8]   (quad) at (6.5,2.0)  {Quadrilatero\\\footnotesize 4 angoli, One Euro};
  \node[box, fill=orange!12] (gest) at (6.5,0)    {Pose della mano\\\footnotesize pugno $\to$ V\\(evento)};
  \node[box, fill=purple!8]  (mask) at (6.5,-2.0) {Maschera\\persona\\\footnotesize $S_t$ e differenza da $B$};
  \node[box, fill=yellow!15, minimum height=18mm] (gpu) at (9.9,0) {Shader GPU\\\footnotesize SDF, filtri,\\ compositing};
  \node[box, fill=gray!10]   (out)  at (13.0,0)   {Schermo,\\ foto, video};
  \draw[arr] (cam) -- (hand);
  \draw[arr] (cam) -- (seg);
  \draw[arr] (hand) -- (quad);
  \draw[arr] (hand) -- (gest);
  \draw[arr] (seg) -- (mask);
  \draw[arr] (quad.east) -- (gpu.north west);
  \draw[arr] (gest) -- (gpu);
  \draw[arr] (mask.east) -- (gpu.south west);
  \draw[arr] (gpu) -- (out);
\end{tikzpicture}
\caption{Flusso di un fotogramma. La segmentazione viene eseguita solo mentre l'effetto
``sparizione'' è attivo, per risparmiare calcolo.}
\label{fig:pipeline}
\end{figure}

Il ciclo principale è guidato da \texttt{requestAnimationFrame}. Il riconoscimento delle mani
viene eseguito solo quando il video ha davvero un fotogramma nuovo: sugli schermi a 60 o 120\,Hz
si evita di analizzare due volte la stessa immagine.

%=====================================================================
\section{Il fotogramma e i sistemi di coordinate}
%=====================================================================

Un fotogramma è una funzione $I:\Omega\to[0,1]^3$ che associa a ogni pixel un colore RGB. Le reti
di MediaPipe restituiscono punti in \emph{coordinate normalizzate} $(u,v)\in[0,1]^2$ rispetto al
fotogramma originale di larghezza $w$ e altezza $h$, con $v$ che cresce verso il basso.

Lo schermo (larghezza $W$, altezza $H$) ha in genere un'altra proporzione, quindi il video viene
scalato ``a copertura'' (\emph{cover}), cioè riempie lo schermo e taglia ciò che avanza:
\begin{equation}
  s=\max\!\left(\frac{W}{w},\frac{H}{h}\right),\qquad
  \vect{o}=\tfrac12\bigl(W-sw,\;H-sh\bigr),\qquad
  \vect{p}=\vect{o}+s\,(w\,u',\,h\,v),
\end{equation}
dove $u'=1-u$ se si usa la fotocamera frontale (effetto specchio, come in un selfie) e $u'=u$
altrimenti. Nello shader si usa la trasformazione inversa: per ogni pixel $\vect{p}$ dello
schermo si calcola quale punto $(u,v)$ del video campionare.

%=====================================================================
\section{Trovare le mani: MediaPipe Hands}
%=====================================================================

Per le mani l'app usa \emph{MediaPipe Hand Landmarker}~\cite{zhang2020hands,lugaresi2019mediapipe},
una pipeline di due reti neurali convoluzionali progettata per girare in tempo reale anche sui
telefoni.

\subsection{Primo stadio: rilevare il palmo}
Una mano è un oggetto articolato: le dita cambiano completamente la sua forma e il suo riquadro.
Il palmo, invece, è quasi rigido. Per questo il primo stadio (\emph{BlazePalm}) non cerca la mano
intera ma il palmo (o il pugno). È un rilevatore a colpo singolo nello stile di
SSD~\cite{liu2016ssd} e BlazeFace~\cite{bazarevsky2019blazeface}. Una rete
\emph{encoder--decoder} produce mappe di caratteristiche a più scale. Per ogni
\emph{ancora} (un riquadro di riferimento) la rete stima la probabilità che contenga un palmo
e la correzione da applicare al riquadro, insieme ad alcuni punti chiave. Siccome i palmi sono
circa quadrati, bastano ancore quadrate, e questo riduce molto il numero di ancore. Per
l'addestramento si usa la \emph{focal loss}~\cite{lin2017focal}, che pesa di più gli esempi
difficili. Infine la \emph{non-maximum suppression} elimina i riquadri duplicati.

Dal palmo trovato si ricava un riquadro \emph{ruotato}: la rotazione allinea il segmento tra
polso e base del dito medio con l'asse verticale. Il riquadro viene poi allargato abbastanza da
contenere le dita.

\subsection{Secondo stadio: i 21 punti}
Il ritaglio ruotato (224$\times$224 pixel) passa a una seconda rete che \emph{regredisce}
direttamente le coordinate di 21 punti anatomici (figura~\ref{fig:hand}), più un punteggio di
presenza della mano e la sua lateralità (destra o sinistra). Ogni punto ha coordinate $(x,y,z)$:
$x$ e $y$ sono normalizzate all'immagine, mentre $z$ è una profondità relativa con origine nel
polso, nella stessa scala di $x$ (valori minori indicano punti più vicini alla fotocamera). La
rete è addestrata su decine di migliaia di foto reali annotate a mano e su immagini sintetiche
di un modello 3D della mano, in pose e sfondi molto vari.

\begin{figure}[htbp]
\centering
\begin{tikzpicture}[scale=0.95,
  pt/.style={circle, fill=dot, inner sep=1.7pt},
  tip/.style={circle, draw=black, thick, fill=dot, inner sep=2.6pt},
  num/.style={font=\scriptsize}]
  \foreach \i/\x/\y in {0/0/0, 1/-0.9/0.5, 2/-1.6/1.1, 3/-2.1/1.7, 4/-2.5/2.25,
     5/-0.8/2.3, 6/-1.0/3.3, 7/-1.1/3.9, 8/-1.2/4.5,
     9/0/2.5, 10/0/3.6, 11/0/4.3, 12/0/4.9,
     13/0.75/2.3, 14/0.85/3.3, 15/0.9/3.9, 16/0.95/4.4,
     17/1.4/1.9, 18/1.65/2.7, 19/1.8/3.2, 20/1.95/3.65}
     \coordinate (p\i) at (\x,\y);
  \foreach \a/\b in {0/1,1/2,2/3,3/4,0/5,5/6,6/7,7/8,5/9,9/10,10/11,11/12,
                     9/13,13/14,14/15,15/16,13/17,17/18,18/19,19/20,0/17}
     \draw[hand, line width=1.6pt] (p\a) -- (p\b);
  \foreach \i in {0,1,2,3,5,6,7,9,10,11,12,13,14,15,16,17,18,19,20} \node[pt] at (p\i) {};
  \node[tip] at (p4) {}; \node[tip] at (p8) {};
  \foreach \i/\d in {0/below, 1/left, 2/left, 3/left, 4/left, 5/left, 6/left, 7/left, 8/above,
     9/right, 10/right, 11/right, 12/above, 13/right, 14/right, 15/right, 16/above,
     17/right, 18/right, 19/right, 20/above}
     \node[num, \d=2pt] at (p\i) {\i};
  \node[font=\small, align=left, anchor=west] at (3.2,3.4) {%
    \textbf{0}\ \ polso\\
    \textbf{1--4}\ \ pollice (CMC, MCP, IP, punta)\\
    \textbf{5--8}\ \ indice (MCP, PIP, DIP, punta)\\
    \textbf{9--12}\ \ medio\\
    \textbf{13--16}\ \ anulare\\
    \textbf{17--20}\ \ mignolo\\[4pt]
    Cerchiati: punte di pollice (4)\\ e indice (8), cioè gli angoli\\ del rettangolo.};
\end{tikzpicture}
\caption{I 21 punti della mano restituiti da MediaPipe e i collegamenti che l'app disegna
(linee rosse, punti verdi, come nel reel di riferimento).}
\label{fig:hand}
\end{figure}

\subsection{Tracciamento nel tempo}
In modalità \texttt{VIDEO} la pipeline non riesegue ogni volta il rilevatore di palmi, che è il
passo più costoso. Il riquadro della mano al fotogramma $t+1$ viene ricavato dai punti trovati
al fotogramma $t$, perché la mano si sposta poco tra due fotogrammi consecutivi. Il rilevatore
torna in gioco solo quando il punteggio di presenza scende sotto la soglia (mano persa) o quando
ci sono meno mani del massimo richiesto (qui 2). Per questo a ogni chiamata va passato un
\emph{timestamp} strettamente crescente.

\paragraph{GPU o CPU.} La stessa rete può essere eseguita con il \emph{delegate} GPU (WebGL) o
CPU (WebAssembly con istruzioni SIMD tramite XNNPACK). L'app prova prima la GPU e passa alla
CPU in caso di errore. L'utente può anche forzare la CPU dalle impostazioni.

\paragraph{Lateralità.} MediaPipe assume che l'immagine in ingresso sia specchiata, come un
selfie. L'app invece passa il fotogramma non specchiato, quindi le etichette ``destra'' e
``sinistra'' risultano scambiate. Per questo vengono usate solo come identificatore stabile di
ciascuna mano (\S\ref{sec:gesti}) e mai mostrate all'utente.

%=====================================================================
\section{Dai punti al quadrilatero}
%=====================================================================

Siano $\vect{a}_k$ e $\vect{b}_k$ i punti $k$ delle due mani, già convertiti in pixel dello
schermo. Le mani vengono ordinate in base alla posizione orizzontale del polso: $\vect{a}$ è
quella più a sinistra sullo schermo.

\begin{description}
  \item[Due mani.] Come nel reel, il quadrilatero è
  $Q=(\vect{a}_8,\,\vect{b}_8,\,\vect{b}_4,\,\vect{a}_4)$: gli indici sono gli angoli in alto e i
  pollici quelli in basso.
  \item[Una mano.] Con punta del pollice $\vect{t}=\vect{p}_4$ e punta dell'indice
  $\vect{i}=\vect{p}_8$, il rettangolo è quello allineato agli assi che ha $\vect{t}$ e $\vect{i}$
  come vertici opposti: $[\min(t_x,i_x),\max(t_x,i_x)]\times[\min(t_y,i_y),\max(t_y,i_y)]$.
\end{description}

Le mani che stanno facendo il gesto pugno o V (\S\ref{sec:gesti}) vengono escluse, così il
gesto non disegna rettangoli per sbaglio.

\paragraph{Quando mostrarlo.} Se le dita si chiudono il quadrilatero degenera. L'area viene
calcolata come somma delle aree (in valore assoluto) dei due triangoli $ABC$ e $ACD$:
\begin{equation}
  \mathcal{A}(Q)=\tfrac12\Bigl(\bigl|(B-A)\times(C-A)\bigr|+\bigl|(C-A)\times(D-A)\bigr|\Bigr),
\end{equation}
dove $\times$ indica il prodotto vettoriale 2D, $(x_1,y_1)\times(x_2,y_2)=x_1y_2-y_1x_2$. A
differenza della formula di Gauss (\emph{shoelace}), questa somma non si annulla quando il
quadrilatero si incrocia ``a farfalla''. L'effetto si mostra solo se
$\mathcal{A}(Q)>0{,}0015\,WH$, cioè se occupa più dello 0,15\% dello schermo.

L'opacità $\alpha$ dell'effetto segue un'esponenziale con costante di tempo $\tau=90$\,ms, e
il quadrilatero resta visibile per 160\,ms dopo l'ultima rilevazione valida:
\begin{equation}
  \alpha\leftarrow\alpha+(\alpha^\star-\alpha)\bigl(1-e^{-\Delta t/\tau}\bigr),\qquad
  \alpha^\star\in\{0,1\}.
\end{equation}
In questo modo l'effetto entra ed esce in dissolvenza e non sfarfalla se la mano viene persa per
un solo fotogramma.

%=====================================================================
\section{Stabilizzare gli angoli: il filtro One Euro}
%=====================================================================

I punti stimati da una rete tremano di qualche pixel da un fotogramma all'altro (\emph{jitter}).
Un filtro passa-basso li calma, ma introduce ritardo (\emph{lag}): il rettangolo ``insegue'' le
dita. Il filtro One Euro~\cite{casiez2012oneeuro} risolve il compromesso adattando la frequenza di
taglio alla velocità del movimento. Per ognuna delle 8 coordinate (4 angoli $\times$ 2 assi):
\begin{align}
  \hat{x}_i &= a_i\,x_i+(1-a_i)\,\hat{x}_{i-1},\qquad
  a_i=\frac{1}{1+\tau_i/\Delta t},\qquad \tau_i=\frac{1}{2\pi f_{c,i}},\\
  f_{c,i} &= f_{\min}+\beta\,\bigl|\hat{\dot{x}}_i\bigr|,
\end{align}
dove $\hat{\dot{x}}_i$ è la velocità stimata (anch'essa filtrata, con taglio fisso a 1\,Hz).

Quando la mano è ferma la velocità è quasi nulla, il taglio resta basso ($f_{\min}$) e il tremolio
viene eliminato. Quando la mano si muove in fretta, $f_c$ sale e il filtro ``si apre'', azzerando
il ritardo. Il cursore \emph{Stabilizzazione} $s\in\{0,\dots,10\}$ dell'app imposta
\begin{equation}
  f_{\min}=\frac{6}{1+0{,}9\,s}\ \text{Hz},\qquad \beta=0{,}004+0{,}0012\,(10-s)\ \text{px}^{-1}.
\end{equation}
Con $s=0$ il filtro è di fatto disattivato. Quando il quadrilatero ricompare dopo essere sparito,
il filtro viene azzerato, così gli angoli non scivolano dalla posizione vecchia.

%=====================================================================
\section{Dal quadrilatero alla maschera: distanza con segno}
\label{sec:sdf}
%=====================================================================

Lo shader deve decidere, per ciascuno dei milioni di pixel dello schermo, se il pixel sta dentro
il quadrilatero e quanto è vicino al bordo. Lo fa con una \emph{funzione di distanza con segno}
(SDF)~\cite{quilez}, negativa dentro e positiva fuori.

\paragraph{Distanza dai lati.} Per un segmento $\overline{\vect{a}\vect{b}}$ si proietta il
pixel $\vect{p}$ sulla retta e si limita la proiezione al segmento:
\begin{equation}
  h=\clamp\!\left(\frac{(\vect{p}-\vect{a})\cdot(\vect{b}-\vect{a})}{\|\vect{b}-\vect{a}\|^2},\,0,\,1\right),
  \qquad d_{ab}(\vect{p})=\bigl\|\vect{p}-\vect{a}-h(\vect{b}-\vect{a})\bigr\|.
\end{equation}
La distanza dal quadrilatero è il minimo sui quattro lati: $d(\vect{p})=\min d_{ab}(\vect{p})$.

\paragraph{Il segno: regola pari-dispari.} Si immagina un raggio orizzontale che parte da
$\vect{p}$ e si contano i lati che attraversa. Un numero dispari di attraversamenti significa
che il pixel è dentro. Il segno è quindi $\sigma=(-1)^{n}$ e $\mathrm{sd}(\vect{p})=-\sigma\,d(\vect{p})$.
Questa regola funziona anche quando i lati si incrociano: un quadrilatero ``a farfalla''
diventa semplicemente due triangoli pieni (figura~\ref{fig:quad}).

\begin{figure}[htbp]
\centering
\begin{tikzpicture}[scale=0.95, lab/.style={font=\scriptsize}]
  \fill[accent!35, even odd rule] (0,2) -- (3,2.2) -- (2.9,0.2) -- (0.1,0) -- cycle;
  \draw[accent!70!black, thick] (0,2) -- (3,2.2) -- (2.9,0.2) -- (0.1,0) -- cycle;
  \node[lab, above left] at (0,2) {$\vect{a}_8$};
  \node[lab, above right] at (3,2.2) {$\vect{b}_8$};
  \node[lab, below right] at (2.9,0.2) {$\vect{b}_4$};
  \node[lab, below left] at (0.1,0) {$\vect{a}_4$};
  \node[font=\small] at (1.5,-0.7) {mani ``a cornice''};
  \begin{scope}[xshift=5.2cm]
  \fill[accent!35, even odd rule] (0,2) -- (3,0.1) -- (3,2.1) -- (0,0) -- cycle;
  \draw[accent!70!black, thick] (0,2) -- (3,0.1) -- (3,2.1) -- (0,0) -- cycle;
  \node[lab, above left] at (0,2) {$\vect{a}_8$};
  \node[lab, below right] at (3,0.1) {$\vect{b}_8$};
  \node[lab, above right] at (3,2.1) {$\vect{b}_4$};
  \node[lab, below left] at (0,0) {$\vect{a}_4$};
  \node[font=\small] at (1.5,-0.7) {mani incrociate: ``farfalla''};
  \end{scope}
  \begin{scope}[xshift=10.4cm]
  \draw[gray] (0,0) rectangle (3.2,2.4);
  \fill[accent!35] (0.6,0.4) rectangle (2.5,1.8);
  \draw[accent!70!black, thick] (0.6,0.4) rectangle (2.5,1.8);
  \node[circle, fill=dot, inner sep=1.5pt, label={[lab]left:$\vect{p}_4$}] at (0.6,0.4) {};
  \node[circle, fill=dot, inner sep=1.5pt, label={[lab]above right:$\vect{p}_8$}] at (2.5,1.8) {};
  \node[font=\small] at (1.5,-0.7) {una mano};
  \end{scope}
\end{tikzpicture}
\caption{Le tre forme possibili. Con la regola pari-dispari il caso incrociato non richiede
nessun trattamento speciale.}
\label{fig:quad}
\end{figure}

\paragraph{Bordi morbidi.} Con la funzione \emph{smoothstep}
\begin{equation}
  \Sm(e_0,e_1,x)=3t^2-2t^3,\qquad t=\clamp\!\left(\frac{x-e_0}{e_1-e_0},0,1\right),
\end{equation}
la maschera vale $m(\vect{p})=1-\Sm(-f,\,f,\,\mathrm{sd}(\vect{p}))$, dove $f$ è la sfumatura
scelta dall'utente più 0,75\,px. Quei 0,75\,px fanno da \emph{anti-aliasing}: il bordo non è mai
seghettato, nemmeno con sfumatura zero. Se il filtro va applicato ``fuori'' si usa $1-m$. Il
colore finale è un'interpolazione lineare, cioè il caso più semplice del compositing
alfa~\cite{porter1984compositing}:
\begin{equation}
  O(\vect{p})=\bigl(1-\alpha\,m(\vect{p})\bigr)\,I(\vect{p})+\alpha\,m(\vect{p})\,F[I](\vect{p}),
\end{equation}
dove $F[I]$ è l'immagine filtrata. Il bordo opzionale è una striscia dove
$|\mathrm{sd}(\vect{p})|$ è minore di metà spessore.

\paragraph{Perché sulla GPU.} Il calcolo di ogni pixel è indipendente dagli altri: è il caso
ideale per un \emph{fragment shader}. Un iPad esegue queste operazioni su circa 2--4 milioni di
pixel per fotogramma senza impegnare la CPU, che resta libera per le reti neurali.

%=====================================================================
\section{I filtri come operatori sull'immagine}
%=====================================================================

Ogni filtro è una funzione $F[I]$ calcolata nello shader. Alcuni sono puntuali (dipendono solo dal
colore del pixel), altri usano un intorno del pixel (convoluzioni) o deformano le coordinate.

\paragraph{Luminanza.} Molti filtri partono dalla luminanza, che pesa i canali in base alla
sensibilità dell'occhio (raccomandazione ITU-R BT.601~\cite{bt601}):
$Y=0{,}299\,R+0{,}587\,G+0{,}114\,B$.

\paragraph{Gradiente di Sobel.} I contorni sono i punti in cui la luminanza cambia
bruscamente. L'operatore di Sobel~\cite{sobel1968,gonzalez2018} stima le derivate con due
convoluzioni $3\times3$:
\begin{equation}
  G_x=\begin{bmatrix}-1&0&1\\-2&0&2\\-1&0&1\end{bmatrix}*Y,\qquad
  G_y=\begin{bmatrix}-1&-2&-1\\0&0&0\\1&2&1\end{bmatrix}*Y,\qquad
  e=\sqrt{G_x^2+G_y^2},\quad \theta=\operatorname{atan2}(G_y,G_x).
\end{equation}

\begin{table}[htbp]
\centering
\small
\begin{tabular}{@{}l>{\raggedright\arraybackslash}p{11.2cm}@{}}
\toprule
\textbf{Filtro} & \textbf{Che cosa calcola} \\
\midrule
Acido & Mappa luminanza e contorni su un arcobaleno che scorre nel tempo:
  $h=\operatorname{frac}(2{,}6\,Y+1{,}4\,e+0{,}1\,n-0{,}12\,t)$, colore $=\mathrm{HSV}(h,\,0{,}85,\,v)$,
  dove $n$ è un rumore pseudo-casuale (grana) e $t$ il tempo. È il look del reel. \\
Bianco e nero & $Y'=\clamp(1{,}15\,(Y-0{,}5)+0{,}5)$: scala di grigi con un leggero aumento di contrasto. \\
Invisibile & Al posto del video mostra lo sfondo catturato $B$ (vedi \S\ref{sec:sparire}). \\
Termico & \emph{Falsi colori}: la luminanza (leggermente sfocata) indicizza una tavolozza lineare a
  tratti nero $\to$ viola $\to$ rosso $\to$ arancio $\to$ bianco, come una termocamera. \\
Negativo & $F[I]=1-I$. \\
Pixel & Quantizza lo spazio: $\vect{q}=(\lfloor\vect{p}/b\rfloor+\tfrac12)\,b$ con blocchi di
  $b=22$ px (riferiti a 720 p), e quantizza i colori a 6 livelli. È un sottocampionamento senza
  filtro anti-alias: l'\emph{aliasing} diventa un effetto estetico. \\
Neon & Contorni su nero, $\Sm(0{,}08,\,0{,}5,\,e)$, con tinta data dall'orientazione $\theta$
  del gradiente: ogni direzione ha il suo colore. \\
Glitch & Aberrazione cromatica: $R(u+\delta)$, $G(u)$, $B(u-\delta)$ con $\delta$ che oscilla, più
  strappi orizzontali casuali a bande e linee di scansione. \\
Fumetto & Posterizzazione a 4 livelli, saturazione aumentata, contorni neri da Sobel e un
  \emph{retino} a punti il cui raggio cresce con $1-Y$ (stampa a mezzatinta). \\
Sfocato & Media di 25 campioni disposti a spirale di Vogel (\S\ref{sec:blur}). \\
Raggi X & Luminanza invertita con correzione gamma, tinta azzurra e contorni luminosi. \\
Caleidoscopio & Coordinate polari attorno al centro del quadrilatero e ``ripiegamento''
  dell'angolo: simmetria a 8 settori. \\
\bottomrule
\end{tabular}
\caption{I 12 filtri e l'operazione di elaborazione d'immagine che li realizza.}
\end{table}

\subsection{Sfocatura con la spirale di Vogel}
\label{sec:blur}
Una sfocatura gaussiana esatta richiederebbe molti campioni. Lo shader usa invece $N=24$ campioni
(più il centro) distribuiti in modo uniforme in un disco con la \emph{spirale di
Vogel}~\cite{vogel1979}, la stessa disposizione dei semi del girasole:
\begin{equation}
  r_k=R\sqrt{k/N},\qquad \varphi_k=k\,\varphi_g,\qquad \varphi_g=\pi\bigl(3-\sqrt5\bigr)\approx2{,}39996,
\end{equation}
dove $\varphi_g$ è l'\emph{angolo aureo}. La radice quadrata rende uniforme la densità sull'area e
l'angolo aureo evita che i campioni si allineino. Il risultato è simile a un \emph{bokeh} circolare
di raggio $R$. La stessa funzione serve alla transizione della sparizione, con $R$ variabile.

\subsection{Caleidoscopio}
Con $\vect{c}$ il centro del quadrilatero, $\vect{d}=\vect{p}-\vect{c}$, $r=\|\vect{d}\|$ e
$\phi=\operatorname{atan2}(d_y,d_x)+\omega t$, l'angolo viene ripiegato in un settore di ampiezza
$\Delta=2\pi/8$:
\begin{equation}
  \phi'=\Bigl|\bigl(\phi\bmod\Delta\bigr)-\tfrac{\Delta}{2}\Bigr|,\qquad
  \vect{p}'=\vect{c}+r\,(\cos\phi',\,\sin\phi').
\end{equation}
Campionare il video in $\vect{p}'$ produce un'immagine con la simmetria del gruppo diedrale: 8
spicchi, ciascuno speculare al vicino.

%=====================================================================
\section{Riconoscere i gesti: pugno e V}
\label{sec:gesti}
%=====================================================================

Per far sparire la persona l'app riconosce una \emph{sequenza} di due pose della stessa mano: un
pugno e poi una V con indice e medio. La classificazione usa regole geometriche sui 21 punti, non
una rete dedicata.

\subsection{Le caratteristiche di ciascun dito}
I punti vengono portati in uno spazio isotropo, $\vect{X}=(x\,w,\;y\,h,\;z\,w)$, in modo che le
distanze abbiano la stessa unità su tutti gli assi. Per ogni dito (escluso il pollice), con
$\vect{M}$ la nocca (MCP), $\vect{P}$ l'articolazione intermedia (PIP), $\vect{T}$ la punta e
$\vect{W}$ il polso, si calcolano due misure:
\begin{align}
  \sigma &= \cos\angle(\vect{P}-\vect{M},\;\vect{T}-\vect{P})
         =\frac{(\vect{P}-\vect{M})\cdot(\vect{T}-\vect{P})}{\|\vect{P}-\vect{M}\|\,\|\vect{T}-\vect{P}\|}
         &&\text{(rettilineità)},\\
  \rho &= \frac{\|\vect{T}-\vect{W}\|}{\|\vect{P}-\vect{W}\|} &&\text{(allungamento)}.
\end{align}
Un dito disteso ha $\sigma\approx1$ (le falangi sono allineate) e $\rho\approx1{,}3$ (la punta è
più lontana dal polso della nocca intermedia). In un dito chiuso a pugno la punta torna verso il
palmo: $\sigma<0$ e $\rho<1$. Entrambe le misure sono rapporti o angoli, quindi non cambiano con
la distanza dalla fotocamera, con la posizione o con la rotazione della mano. Siccome usano anche
$z$, funzionano pure quando le dita puntano verso l'obiettivo.

\subsection{Le regole}
\begin{align*}
  \textsf{disteso}   &\iff \sigma>0{,}55\ \wedge\ \rho>1{,}08,\\
  \textsf{piegato}   &\iff \sigma<0{,}25\ \vee\ \rho<0{,}92,\\
  \textsf{pugno}     &\iff \text{indice, medio, anulare, mignolo tutti \textsf{piegati}},\\
  \textsf{V}         &\iff \text{indice e medio \textsf{distesi}}\ \wedge\ \text{anulare e mignolo \textsf{piegati}}
                        \ \wedge\ \frac{\|\vect{T}_8-\vect{T}_{12}\|}{\|\vect{X}_9-\vect{W}\|}>0{,}3.
\end{align*}
L'ultima condizione richiede che le due dita siano \emph{divaricate}: due dita unite non sono una
V. Il pollice viene ignorato, perché in entrambe le pose può stare ovunque. Le regole sono
volutamente asimmetriche: per ``disteso'' servono entrambe le condizioni, per ``piegato'' ne basta
una. Così una V viene riconosciuta solo quando è netta, mentre un pugno un po' lasco conta lo
stesso.

\subsection{La sequenza nel tempo}
Le pose stimate fotogramma per fotogramma sono rumorose, quindi l'app usa una piccola macchina a
stati per ogni mano, identificata dalla sua lateralità:
\begin{enumerate}
  \item \textbf{Anti-rimbalzo} (\emph{debounce}): una posa diventa ``ufficiale'' solo se resta
  uguale per almeno 70\,ms.
  \item \textbf{Pugno valido}: il pugno deve durare almeno 220\,ms. Da quel momento, e per tutto
  il tempo in cui resta chiuso, l'app aggiorna l'istante $t_{\text{pugno}}$. Sullo schermo
  compare ``ora V''.
  \item \textbf{V}: la V deve iniziare entro 1{,}6\,s dall'ultimo pugno valido e durare almeno
  180\,ms. A quel punto scatta l'evento.
  \item \textbf{Pausa} (\emph{cooldown}): per 1{,}2\,s non può scattare un altro evento.
\end{enumerate}

\begin{figure}[htbp]
\centering
\begin{tikzpicture}[x=4.2cm, y=1cm, font=\small]
  \draw[-{Stealth}] (0,0) -- (3.05,0) node[right] {$t$ (s)};
  \foreach \t in {0,0.5,1,1.5,2,2.5,3} \draw (\t,0.06) -- (\t,-0.06) node[below, font=\scriptsize] {\t};
  \fill[gray!20]   (0,0.3) rectangle (0.3,0.9);  \node at (0.15,0.6) {\scriptsize altro};
  \fill[orange!45] (0.3,0.3) rectangle (1.2,0.9); \node at (0.75,0.6) {pugno};
  \fill[gray!20]   (1.2,0.3) rectangle (1.5,0.9); \node at (1.35,0.6) {\scriptsize altro};
  \fill[accent!45] (1.5,0.3) rectangle (2.4,0.9); \node at (1.95,0.6) {V};
  \fill[gray!20]   (2.4,0.3) rectangle (3,0.9);  \node at (2.7,0.6) {\scriptsize altro};
  \draw[decorate, decoration={brace, amplitude=4pt}] (0.3,1.0) -- node[above=4pt, font=\scriptsize] {220 ms} (0.52,1.0);
  \draw[dashed] (0.52,0.3) -- (0.52,1.0);
  \draw[decorate, decoration={brace, amplitude=4pt}] (1.5,1.0) -- node[above=4pt, xshift=-7pt, font=\scriptsize] {180 ms} (1.68,1.0);
  \draw[decorate, decoration={brace, mirror, amplitude=4pt}] (1.2,-0.55) -- node[below=4pt, font=\scriptsize] {la V deve iniziare entro 1,6 s dall'ultimo pugno valido} (2.8,-0.55);
  \node[star, star points=5, fill=hand, inner sep=2.2pt] (ev) at (1.68,1.55) {};
  \node[right=2pt of ev, font=\scriptsize] {evento: sparisci / riappari};
  \draw[hand, thick] (1.68,0.3) -- (1.68,1.45);
\end{tikzpicture}
\caption{Diagramma temporale della sequenza pugno $\to$ V.}
\end{figure}

Chiedere due pose diverse \emph{in quest'ordine} riduce molto i falsi positivi. Una V da sola,
gesto frequente in foto e video, non fa scattare niente. Ogni evento \emph{inverte} lo stato:
il primo fa sparire, il secondo fa riapparire.

%=====================================================================
\section{Sparire: segmentazione, sottrazione dello sfondo e compositing}
\label{sec:sparire}
%=====================================================================

Per ``togliere'' una persona serve sapere due cose: \emph{quali} pixel le appartengono e
\emph{che cosa} c'era dietro.

\subsection{Che cosa c'era dietro: la lastra di sfondo}
L'app cattura una \emph{lastra pulita} (\emph{clean plate}) $B$, cioè un fotogramma della scena
vuota, durante un conto alla rovescia in cui l'utente esce dall'inquadratura. L'ipotesi è che la
fotocamera resti ferma e la luce costante: in quel caso, in ogni pixel non coperto dalla persona,
vale $I_t(\vect{p})\approx B(\vect{p})$.

\subsection{Quali pixel sono tuoi: due indizi complementari}

\paragraph{1. Segmentazione semantica.} Il \emph{Selfie Segmenter} di MediaPipe è una piccola rete
convoluzionale \emph{encoder--decoder} derivata da MobileNetV3~\cite{howard2019mobilenetv3}.
Riceve l'immagine ridotta a 256$\times$256 e restituisce per ogni pixel la probabilità
$S(\vect{p})\in[0,1]$ che appartenga a una persona. La maschera viene poi riportata alla
risoluzione del fotogramma. È una stima ``semantica'': riconosce una persona anche se ha i
vestiti dello stesso colore del muro, ma tende a sbagliare sui contorni (capelli, dita) e ignora
le ombre e gli oggetti che la persona tiene in mano.

\paragraph{2. Sottrazione dello sfondo.} Il secondo indizio è il più classico della
videosorveglianza~\cite{piccardi2004background}: ciò che è diverso dalla lastra vuota è primo
piano. Per ridurre il rumore del sensore si confrontano medie locali su 5 pixel a croce:
\begin{equation}
  D(\vect{p})=\bigl\|\bar I(\vect{p})-\bar B(\vect{p})\bigr\|,\qquad
  M_d(\vect{p})=\Sm(0{,}09,\,0{,}2,\,D(\vect{p})).
\end{equation}
Questo indizio coglie ombre, oggetti tenuti in mano e i pezzi che la rete non vede. Fallisce però
dove la persona ha lo stesso colore dello sfondo (\emph{camouflage}). I punti deboli dei due
metodi sono diversi, quindi conviene combinarli.

\subsection{Costi asimmetrici: meglio coprire troppo}
L'osservazione chiave è che i due errori possibili non costano uguale:
\begin{itemize}
  \item \textbf{falso negativo} (un pixel della persona non viene sostituito): resta un ``fantasma''
  ben visibile;
  \item \textbf{falso positivo} (un pixel dello sfondo viene sostituito): non si vede, perché lì
  $I\approx B$ e si sostituisce lo sfondo con se stesso.
\end{itemize}
Per questo tutte le scelte spingono la maschera ad \emph{allargarsi}:
\begin{align}
  S_d(\vect{p}) &= \max_{\vect{q}\in\mathcal{N}(\vect{p})} S(\vect{q})
     &&\text{dilatazione morfologica~\cite{serra1982} su 9 campioni, raggio 14 px},\\
  M_s(\vect{p}) &= \Sm(0{,}05,\,0{,}3,\,S_d(\vect{p}))
     &&\text{soglia bassa e morbida},\\
  M(\vect{p}) &= \max\bigl(M_s(\vect{p}),\,M_d(\vect{p})\bigr)
     &&\text{unione (``OR'' sfumato) dei due indizi}.
\end{align}
Anche nel tempo vale lo stesso principio. La maschera della rete viene tenuta con un
\emph{peak hold} che decade lentamente,
\begin{equation}
  S_t(\vect{p})=\max\bigl(\hat S_t(\vect{p}),\;0{,}8\,S_{t-1}(\vect{p})\bigr),
\end{equation}
così un pixel appena lasciato dalla persona resta coperto per qualche fotogramma e i bordi non
sfarfallano.

\subsection{La transizione: sfocatura e dissolvenza}
Un parametro $t\in[0,1]$ va da 0 a 1 in 0{,}6\,s (o al contrario per riapparire). Il colore
``fantasma'' $G$ e l'uscita $O$ sono
\begin{align}
  r(t) &= 26\,\Sm(0,\,0{,}55,\,t)\ \text{px}, \qquad w(t)=\Sm(0{,}3,\,1,\,t),\\
  G(\vect{p}) &= \bigl(1-w(t)\bigr)\,\mathrm{Blur}_{r(t)}[I](\vect{p})+w(t)\,B(\vect{p}),\\
  O(\vect{p}) &= \bigl(1-M(\vect{p})\bigr)\,I(\vect{p})+M(\vect{p})\,G(\vect{p}).
\end{align}
Nella prima metà la persona si sfoca rapidamente ($r$ raggiunge quasi il massimo a metà
transizione), poi si dissolve nello sfondo ($w$ cresce da 0,3 in poi). A $t=1$ resta solo $B$
dentro la maschera, mentre fuori la scena è ancora dal vivo.

\begin{figure}[htbp]
\centering
\begin{tikzpicture}[x=7cm, y=3.2cm,
  declare function={sst(\a,\b,\x)=3*(min(max((\x-\a)/(\b-\a),0),1))^2-2*(min(max((\x-\a)/(\b-\a),0),1))^3;}]
  \draw[-{Stealth}] (0,0) -- (1.08,0) node[right] {$t$};
  \draw[-{Stealth}] (0,0) -- (0,1.22);
  \foreach \t in {0,0.25,0.5,0.75,1} \draw (\t,0.03) -- (\t,-0.03) node[below, font=\scriptsize] {\t};
  \draw (0.015,1) -- (-0.015,1) node[left, font=\scriptsize] {1};
  \draw[very thick, blue!70!black, domain=0:1, samples=80] plot (\x, {sst(0,0.55,\x)});
  \draw[very thick, accent, domain=0:1, samples=80] plot (\x, {sst(0.3,1,\x)});
  \node[blue!70!black, font=\small, anchor=west] at (0.6,1.1) {$r(t)/26$: sfocatura};
  \node[accent!80!black, font=\small, anchor=west] at (0.62,0.35) {$w(t)$: dissolvenza};
\end{tikzpicture}
\caption{Le due curve della transizione.}
\end{figure}

\begin{table}[htbp]
\centering
\small
\begin{tabular}{@{}cccl@{}}
\toprule
$t$ & $r(t)$ & $w(t)$ & dentro la maschera si vede\\
\midrule
0    & 0 px    & 0    & la persona nitida\\
0,25 & 11,2 px & 0    & la persona già sfocata\\
0,5  & 25,4 px & 0,20 & una sagoma molto sfocata che inizia a svanire\\
0,75 & 26 px   & 0,71 & un fantasma trasparente\\
1    & ---     & 1    & solo lo sfondo $B$\\
\bottomrule
\end{tabular}
\caption{Valori della transizione nei suoi punti chiave.}
\end{table}

Il cuore dello shader è breve:
\begin{lstlisting}
float m = personMask(uv);                  // M: segmentation OR difference
vec3 bg = texture2D(uBg, uv).rgb;          // B: clean plate
float radius = 26.0 * smoothstep(0.0, 0.55, t);
vec3 ghost = mix(blurAt(uv, radius), bg, smoothstep(0.3, 1.0, t));
col = mix(live, ghost, m);                 // O = (1-M) I + M G
\end{lstlisting}

\subsection{Limiti}
\begin{itemize}
  \item \textbf{Fotocamera mobile}: se l'inquadratura si sposta, $B$ non è più allineata. Servirebbe
  registrare i fotogrammi, ad esempio stimando un'omografia con punti caratteristici, oppure
  ricostruire lo sfondo nel tempo.
  \item \textbf{Luce ed esposizione automatica}: se la luminosità cambia, $D$ cresce ovunque. La
  maschera si allarga (senza conseguenze visibili) ma lo sfondo mostrato ha la luce ``vecchia''.
  \item \textbf{Zone mai viste}: la lastra contiene solo ciò che era visibile quando la scena era
  vuota. Per questo va catturata uscendo davvero dall'inquadratura.
  \item \textbf{Altre persone o oggetti in movimento} vengono trattati come primo piano dalla
  sottrazione dello sfondo e quindi nascosti anch'essi.
\end{itemize}

%=====================================================================
\section{Tempo reale e privacy}
%=====================================================================

A 30 fotogrammi al secondo ogni fotogramma ha un budget di circa 33\,ms. Per rispettarlo l'app:
\begin{itemize}
  \item analizza solo i fotogrammi nuovi e sfrutta il tracciamento di MediaPipe, che evita di
  rieseguire il rilevatore di palmi;
  \item esegue la segmentazione solo mentre la sparizione è attiva o in transizione;
  \item affida alla GPU tutto il lavoro per pixel. Il costo dipende soprattutto dal numero di
  letture di texture per pixel: 25 per la sfocatura, 8 per Sobel, una ventina per la maschera
  della persona;
  \item permette di ridurre la risoluzione di rendering dalle impostazioni sui dispositivi meno
  potenti.
\end{itemize}
Nessuna immagine lascia il dispositivo: modelli e librerie vengono scaricati una volta, poi
l'inferenza e il compositing avvengono interamente nel browser.

%=====================================================================
\begin{thebibliography}{99}
%=====================================================================

\bibitem{zhang2020hands}
F.~Zhang, V.~Bazarevsky, A.~Vakunov, A.~Tkachenka, G.~Sung, C.-L.~Chang, M.~Grundmann.
\newblock \emph{MediaPipe Hands: On-device Real-time Hand Tracking}.
\newblock arXiv:2006.10214, 2020.

\bibitem{lugaresi2019mediapipe}
C.~Lugaresi et al.
\newblock \emph{MediaPipe: A Framework for Building Perception Pipelines}.
\newblock arXiv:1906.08172, 2019.

\bibitem{bazarevsky2019blazeface}
V.~Bazarevsky, Y.~Kartynnik, A.~Vakunov, K.~Raveendran, M.~Grundmann.
\newblock \emph{BlazeFace: Sub-millisecond Neural Face Detection on Mobile GPUs}.
\newblock arXiv:1907.05047, 2019.

\bibitem{liu2016ssd}
W.~Liu, D.~Anguelov, D.~Erhan, C.~Szegedy, S.~Reed, C.-Y.~Fu, A.~C.~Berg.
\newblock \emph{SSD: Single Shot MultiBox Detector}.
\newblock ECCV, 2016.

\bibitem{lin2017focal}
T.-Y.~Lin, P.~Goyal, R.~Girshick, K.~He, P.~Doll\'ar.
\newblock \emph{Focal Loss for Dense Object Detection}.
\newblock ICCV, 2017.

\bibitem{howard2019mobilenetv3}
A.~Howard et al.
\newblock \emph{Searching for MobileNetV3}.
\newblock ICCV, 2019.

\bibitem{casiez2012oneeuro}
G.~Casiez, N.~Roussel, D.~Vogel.
\newblock \emph{1\,\texteuro{} Filter: A Simple Speed-based Low-pass Filter for Noisy Input in Interactive Systems}.
\newblock CHI, 2012.

\bibitem{quilez}
I.~Quilez.
\newblock \emph{2D distance functions}.
\newblock \url{https://iquilezles.org/articles/distfunctions2d/}.

\bibitem{porter1984compositing}
T.~Porter, T.~Duff.
\newblock \emph{Compositing Digital Images}.
\newblock SIGGRAPH, 1984.

\bibitem{bt601}
ITU-R.
\newblock \emph{Recommendation BT.601: Studio encoding parameters of digital television}.

\bibitem{sobel1968}
I.~Sobel, G.~Feldman.
\newblock \emph{A 3$\times$3 Isotropic Gradient Operator for Image Processing}.
\newblock Stanford Artificial Intelligence Project, 1968.

\bibitem{gonzalez2018}
R.~C.~Gonzalez, R.~E.~Woods.
\newblock \emph{Digital Image Processing}, 4\textsuperscript{a} ed.
\newblock Pearson, 2018.

\bibitem{vogel1979}
H.~Vogel.
\newblock \emph{A better way to construct the sunflower head}.
\newblock Mathematical Biosciences, 44(3--4), 1979.

\bibitem{piccardi2004background}
M.~Piccardi.
\newblock \emph{Background subtraction techniques: a review}.
\newblock IEEE International Conference on Systems, Man and Cybernetics, 2004.

\bibitem{serra1982}
J.~Serra.
\newblock \emph{Image Analysis and Mathematical Morphology}.
\newblock Academic Press, 1982.

\end{thebibliography}

\end{document}
